% =====================================================================
% The principal W_3 datum: exact local algebra and reconstruction gates.
% Exact-rational companion:
% compute/lib/theorem_w3_holographic_datum_engine.py.
% =====================================================================

\providecommand{\BRST}{\mathrm{BRST}}
\providecommand{\BRSTGauged}{\mathrm{BRSTGaugedChirAlg}}

\chapter{The principal \texorpdfstring{$\mathcal W_3$}{W3} datum:
exact algebra and reconstruction}
\label{ch:w3-holographic-datum}

The principal algebra \(\mathcal W_3^k\) is the first chart in the
standard landscape with two strong generators of different weights.
Its local vertex-algebra structure is explicit.  The passage from this
local structure to a seven-entry holographic modular Koszul datum is a
sequence of reconstruction problems.  Keeping these two layers apart
reveals the precise mathematical content of the example.

The exact layer consists of the Fateev--Lukyanov central charge, the
Zamolodchikov OPE, the norm and zero mode of the Virasoro
quasi-primary \(\Lambda\), and the level-one Shapovalov determinant.
The reconstruction layer consists of a presentation coalgebra,
completed chiral bar data, DS/bar transport, the derived centre,
collision residues, modular traces, stable-graph sewing, scalar
shadows, commuting Hamiltonians, and an open--closed comparison.
Each reconstruction step below carries its own hypothesis package.

\section{The hypothesis ledger}
\label{sec:w3hol-seven-entries}

Fix a smooth projective curve \(X\), a marked divisor \(D\), a
tangential structure \(\tau\), and a principal level
\(k\in\mathbb C\) with \(k+3\in\mathbb C^\times\).  We use the
following packages.
\begin{description}
\item[\(H_{W_3}^{\mathrm{pres}}\).]
A topological presentation
\(\mathcal W_3^k\simeq T_X(V)/(R)\), including the nonlinear
\(\Lambda\)-relation, together with the presentation coalgebra
\[
(\mathcal W_3^k)^{\mathrm i}:=C_X(s^{-1}V,s^{-2}R)
\quad\text{and}\quad
q_{\mathcal W_3}\colon
(\mathcal W_3^k)^{\mathrm i}\longrightarrow
\bar B_X^{\mathrm{ch}}(\mathcal W_3^k).
\]
\item[\(H_{W_3}^{\mathrm{bar}}(X,D,\tau)\).]
An ordered configuration-space bar complex with fixed suspension
signs, conilpotent or complete topology, and continuous twisting
coderivation.
\item[\(H_{W_3}^{\mathrm{DS/bar}}\).]
PBW convergence, BRST concentration, a filtered comparison between
DS reduction and the completed bar construction, and continuous
Ran--Verdier duality.
\item[\(H_{W_3}^{\mathrm{cen}}\).]
A chain model for
\(Z_{\mathrm{ch}}^{\mathrm{der}}(\mathcal W_3^k)\), together with the
family support datum
\(H_H(\mathcal W_3^k;S_{W_3})\).
\item[\(H_{W_3}^{\mathrm{coll}}\).]
A collision-residue map on the ordered bar complex which retains
descendants and composite channels and fixes the local coordinate
and \(d\log E\) convention.
\item[\(H_{W_3}^{\mathrm{mod}}\).]
A trace-compatible genus-one curvature construction defining the
scalar \(\kappa_{\mathrm{ch}}(\mathcal W_3^k)\).
\item[\(H_{W_3}^{\mathrm{res}}\) and
\(H_{W_3}^{\mathrm{proj}}\).]
Chain-compatible residue and scalar-projection maps from the full
multi-channel Maurer--Cartan tensor.
\item[\(H_{W_3}^{\mathrm{rescale}}\).]
A family contraction \((\iota_k,p_k,h_k,\nu_k)\) on one completed
bar--cobar complex and a scalar trace identifying \(\nu_k\) with the
algebra-level Verdier sum \(K^\kappa(\mathcal W_3^k)\).
\item[\(H_{W_3}^{\mathrm{sew}}\).]
Clutching-compatible propagators, automorphism factors, and
convergence for the stable-graph expansion.
\item[\(H_{W_3}^{\mathrm{flat}}\).]
A representation of the ordered Maurer--Cartan element on conformal
blocks for which the induced connection is flat.
\item[\(H_{W_3}^{\mathrm{line}}\) and
\(H_{W_3}^{\mathrm{OC}}\).]
A comparison with the chosen line-operator category and an
open--closed map from the physical HT observable algebra to the
chiral derived centre.
\end{description}
The total holographic package is
\[
H_{W_3}^{\mathrm{hol}}
:=
H_{W_3}^{\mathrm{pres}}+
H_{W_3}^{\mathrm{bar}}+
H_{W_3}^{\mathrm{DS/bar}}+
H_{W_3}^{\mathrm{cen}}+
H_{W_3}^{\mathrm{coll}}+
H_{W_3}^{\mathrm{mod}}+
H_{W_3}^{\mathrm{res}}+
H_{W_3}^{\mathrm{proj}}+
H_{W_3}^{\mathrm{rescale}}+
H_{W_3}^{\mathrm{sew}}+
H_{W_3}^{\mathrm{flat}}+
H_{W_3}^{\mathrm{line}}+
H_{W_3}^{\mathrm{OC}}.
\]

\subsection{The chart algebra \texorpdfstring{$\mathcal A=\mathcal W_3^k$}{A=W3}}
\label{subsec:w3hol-A}

The principal DS algebra
\[
\mathcal A=\mathcal W^k(\mathfrak{sl}_3,f_{\mathrm{prin}})
\]
is strongly generated on the universal generic fibre by a Virasoro
field \(T\) of weight \(2\) and a Virasoro-primary field \(W\) of
weight \(3\).  Its central charge is
\begin{equation}
\label{eq:w3hol-central-charge}
c(k)=2-\frac{24(k+2)^2}{k+3}.
\end{equation}
Put
\[
\Lambda={:}TT{:}-\frac{3}{10}\partial^2T,
\qquad D(c):=5c+22.
\]
In the Zamolodchikov normalization, with \(D(c)\in\mathbb C^\times\),
\begin{align*}
T(z)T(w)&\sim
\frac{c/2}{(z-w)^4}
+\frac{2T(w)}{(z-w)^2}
+\frac{\partial T(w)}{z-w},\\
T(z)W(w)&\sim
\frac{3W(w)}{(z-w)^2}
+\frac{\partial W(w)}{z-w},\\
W(z)W(w)&\sim
\frac{c/3}{(z-w)^6}
+\frac{2T(w)}{(z-w)^4}
+\frac{\partial T(w)}{(z-w)^3}\\
&\quad+
\frac{\frac{3}{10}\partial^2T(w)
+\frac{32}{D(c)}\Lambda(w)}{(z-w)^2}
+\frac{\frac{1}{15}\partial^3T(w)
+\frac{16}{D(c)}\partial\Lambda(w)}{z-w}.
\end{align*}
This is the Zamolodchikov algebra
\cite{Zamolodchikov,Bouwknegt-Schoutens}; the factor \(32/16\)
is the OPE/translation-covariance pair obtained by the double-contour
conversion from singular products to mode commutators.

\subsection{The presentation coalgebra}
\label{subsec:w3hol-A-i}

Under \(H_{W_3}^{\mathrm{pres}}\), the second entry is
\begin{equation}\label{eq:w3hol-ai}
\mathcal A^{\mathrm i}
=C_X(s^{-1}V,s^{-2}R)
\xrightarrow{\ q_{\mathcal A}\ }
\bar B_X^{\mathrm{ch}}(\mathcal A).
\end{equation}
The full bar complex retains its differential and completion.
Koszulity is the assertion that \(H^\bullet(q_{\mathcal A})\) is an
isomorphism on the chosen locus.  Bar--cobar reconstruction then
recovers the chart through
\(\Omega_X^{\mathrm{ch}}\bar B_X^{\mathrm{ch}}(\mathcal A)
\simeq\mathcal A\).
For the nonlinear \(W_3\) presentation, constructing
\((V,R,q_{\mathcal A})\) with these properties is an open finite
presentation problem.

\subsection{The Ran--Verdier and Koszul entries}
\label{subsec:w3hol-A-dual}

The completed Ran--Verdier algebra is
\[
K_X(\mathcal A)
:=\mathbb D_{\operatorname{Ran}}
\bar B_X^{\mathrm{ch}}(\mathcal A).
\]
Under \(H_{W_3}^{\mathrm{DS/bar}}\), a chosen rectification map
\[
\nu_{\mathcal A}\colon
K_X(\mathcal A)\longrightarrow\mathcal A^!
\]
may identify the strict Koszul entry with the principal algebra at
the reflected parameter \(k^\vee=-k-6\).  The exact arithmetic of
this reflected parameter is proved in
Theorem~\ref{thm:w3hol-conductor}; the algebra equivalence has status
\ClaimStatusConditional under \(H_{W_3}^{\mathrm{DS/bar}}\).

\subsection{The derived centre and the line comparison}
\label{subsec:w3hol-C}

The fourth algebraic entry is prescribed by the chiral derived centre
\begin{equation}\label{eq:w3hol-derived-centre}
\mathcal C
:=Z_{\mathrm{ch}}^{\mathrm{der}}(\mathcal W_3^k)
\simeq
C_{\mathrm{ch}}^\bullet(\mathcal W_3^k,\mathcal W_3^k).
\end{equation}
The package \(H_{W_3}^{\mathrm{cen}}\) constructs a family chain model
and proves
\[
\operatorname{Supp}
\ChirHoch^\bullet(\mathcal W_3^k)
\subseteq S_{W_3}.
\]
A perfect degree-$d$ pairing is a separate input for any reflected
degree statement on the Koszul companion.
The evaluation-generated line category has the conditional comparison
\begin{equation}\label{eq:w3hol-line-evaluation-category}
\mathcal C_{\mathrm{line}}^{\mathrm{ev}}
\longrightarrow
\operatorname{Rep}_q(\mathfrak{sl}_3),
\qquad q=\exp\!\left(\frac{\pi i}{k+3}\right),
\end{equation}
under \(H_{W_3}^{\mathrm{line}}\).  The physical comparison consists
of typed maps
\begin{equation}\label{eq:w3hol-typed-comparison-maps}
A_{\mathrm{bulk}}^{\mathrm{HT}}
\xrightarrow{\ \beta_{\mathrm{der}}\ }
Z_{\mathrm{ch}}^{\mathrm{der}}(\mathcal W_3^k),
\qquad
\mathcal C_{\mathrm{line}}^{\mathrm{ev}}
\xrightarrow{\ \iota_{\mathrm{ev}}\ }
\mathcal C_{\mathrm{line}},
\end{equation}
with Swiss-cheese compatibility.  These maps constitute
\(H_{W_3}^{\mathrm{OC}}+H_{W_3}^{\mathrm{line}}\).

\subsection{The collision entry}
\label{subsec:w3hol-r}

Under \(H_{W_3}^{\mathrm{bar}}+H_{W_3}^{\mathrm{coll}}\), the binary
ordered-bar entry is
\begin{equation}\label{eq:w3hol-r-collision}
K_{\mathcal W_3}^{\mathrm{coll}}(z)
:=
\operatorname{Res}_{0,2}^{\mathrm{coll}}
\bigl(\Theta_{\mathcal W_3}^{E_1}\bigr)
\in
\operatorname{End}_{\mathcal W_3}(2)\otimes\mathcal O(*\Delta).
\end{equation}
A represented classical \(r\)-matrix is the image of this kernel
under \(H_{\mathrm{CYBE}}^{\mathrm{rep}}\).  The local OPE determines
the candidate channel coefficients; the residue theorem identifies
them with the ordered-bar kernel.

\subsection{The Maurer--Cartan entry}
\label{subsec:w3hol-Theta}

The completed ordered convolution algebra is
\[
\mathfrak g_{\mathcal W_3}^{E_1}
:=
\operatorname{Conv}\!\left(
B^{\mathrm{ord}}(\mathcal W_3^k),\mathcal W_3^k
\right).
\]
Under \(H_{W_3}^{\mathrm{bar}}\), the twisting coderivation defines
\(\Theta_{\mathcal W_3}\in\mathfrak g_{\mathcal W_3}^{E_1}[1]\)
satisfying
\[
D\Theta_{\mathcal W_3}
+\frac12[\Theta_{\mathcal W_3},\Theta_{\mathcal W_3}]=0.
\]
The presentation comparison and the full multi-channel completion are
the concrete \(W_3\) obligations in this entry.

\subsection{The holomorphic connection entry}
\label{subsec:w3hol-nabla}

Under
\(H_{W_3}^{\mathrm{bar}}+H_{W_3}^{\mathrm{res}}
+H_{W_3}^{\mathrm{flat}}\), the represented shadow map defines
\begin{equation}\label{eq:w3hol-shadow-connection}
\nabla_{0,n}^{\mathrm{hol}}
:=
d-\operatorname{Sh}_{0,n}(\Theta_{\mathcal W_3}).
\end{equation}
The Maurer--Cartan equation supplies the curvature identity after the
representation and convergence hypotheses have been verified.

\section{Reflected central arithmetic}
\label{sec:w3hol-complementarity}

\begin{theorem}[Reflected principal central sum;
\ClaimStatusProvedHere]
\label{thm:w3hol-conductor}
\emph{Type signature:}
\((\mathrm{Open},\mathrm{universal\ principal\ DS},1,
H_{\mathrm{univ}}^{W_3})\).
Let \(k^\vee=-k-6\).  Then
\begin{equation}\label{eq:w3hol-conductor}
c(k)+c(k^\vee)=100.
\end{equation}
The induced affine reflection on the central-charge coordinate is
\(c\mapsto100-c\), with arithmetic midpoint
\begin{equation}\label{eq:w3hol-self-dual-point}
c_{\mathrm{mid}}=50.
\end{equation}
\end{theorem}

\begin{proof}
Set \(t=k+3\).  Formula~\eqref{eq:w3hol-central-charge} becomes
\[
c(t)=2-24\frac{(t-1)^2}{t}.
\]
The reflected level has \(t^\vee=-t\), hence
\[
c(-t)=2+24\frac{(t+1)^2}{t}.
\]
Therefore
\[
c(t)+c(-t)
=4+24\frac{(t+1)^2-(t-1)^2}{t}
=4+96
=100.
\]
The midpoint of the reflected \(c\)-coordinate is \(100/2=50\).
\end{proof}

\begin{remark}[Arithmetic and reconstruction]
\label{rem:w3hol-w-complementarity}
Equation~\eqref{eq:w3hol-conductor} is an identity in
\(\mathbb C(k)\).  The package \(H_{W_3}^{\mathrm{DS/bar}}\) promotes
the reflected parameter to a comparison of completed bar--Verdier
objects.  The KSDual fixed locus is then defined inside the resulting
moduli problem.
\end{remark}

\begin{remark}[Three scalar coordinates]
\label{rem:w3hol-self-dual-vs-critical}
The number \(50\) is the midpoint of the reflected central coordinate.
A BRST critical charge is computed from a specified ghost complex and
nilpotence equation.  A KSDual point is a fixed object of the
Verdier--Koszul involution.  These three constructions live at levels
\(1\), \(5\), and \(1\leftrightarrow2\), respectively.
\end{remark}

\section{The exact composite and highest-weight packet}
\label{sec:w3hol-lambda}

The Virasoro calculation gives
\[
|\Lambda\rangle
=L_{-2}^2|0\rangle-\frac35L_{-4}|0\rangle,
\qquad
N_\Lambda
:=\langle\Lambda|\Lambda\rangle
=\frac{c(5c+22)}{10}.
\]
Thus the one-dimensional inverse Gram factor is
\[
N_\Lambda^{-1}=\frac{10}{c(5c+22)}
\]
on the regular Zamolodchikov surface
\(c(5c+22)\in\mathbb C^\times\).  A scalar quartic shadow additionally
uses \(H_{W_3}^{\mathrm{res}}+H_{W_3}^{\mathrm{proj}}\).

\begin{proposition}[Action of
\texorpdfstring{$\Lambda_0$}{Lambda0} on primaries;
\ClaimStatusProvedHere]
\label{prop:w3hol-lambda-on-primaries}
\emph{Type signature:}
\((\mathrm{Open},\mathrm{Zamolodchikov\ highest\ weight},1,
H_{\mathrm{hw}}^{W_3})\).
For a highest-weight vector \(|h,w\rangle\),
\begin{equation}\label{eq:w3hol-lambda-action}
\Lambda_0|h,w\rangle
=\left(h^2+\frac h5\right)|h,w\rangle.
\end{equation}
\end{proposition}

\begin{proof}
The weight-four mode formula
\[
\Lambda_n
=\sum_{p\in\mathbb Z}{:}L_pL_{n-p}{:}
-\frac3{10}(n+2)(n+3)L_n
\]
gives
\begin{equation}\label{eq:w3hol-tt-zero-mode}
({:}TT{:})_0
=2\sum_{m\geq1}L_{-m}L_m+L_0^2+2L_0.
\end{equation}
On \(|h,w\rangle\), this acts by \(h^2+2h\).  Since
\((\partial^2T)_0=6L_0\), the correction
\(-\frac3{10}(\partial^2T)_0\) acts by \(-9h/5\).  Their sum is
\(h^2+h/5\).
\end{proof}

\begin{corollary}[The conformal roots of
\texorpdfstring{$\Lambda_0$}{Lambda0};
\ClaimStatusProvedHere]
\label{cor:w3hol-lambda-roots}
\emph{Type signature:}
\((\mathrm{Open},\mathrm{Zamolodchikov\ highest\ weight},1,
H_{\mathrm{hw}}^{W_3})\).
The \(\Lambda_0\)-eigenvalue vanishes at
\[
h=0,\qquad h=-\frac15.
\]
\end{corollary}

\begin{remark}[Central-charge dependence]
\label{rem:w3hol-lambda-c-independent}
The highest-weight action of \(\Lambda_0\) is determined by Virasoro
normal ordering.  The central charge enters the \(WW\) product through
\(32/(5c+22)\) and enters the Gram norm through
\(N_\Lambda=c(5c+22)/10\).
\end{remark}

In the ordered basis
\(\{L_{-1}|h,w\rangle,W_{-1}|h,w\rangle\}\), the level-one
Shapovalov matrix is
\begin{equation}\label{eq:w3hol-level-one-gram}
M_1(c,h,w)=
\begin{pmatrix}
2h&3w\\[2pt]
3w&
-\frac h5+\frac{32}{5c+22}
\left(h^2+\frac h5\right)
\end{pmatrix}.
\end{equation}
Its determinant satisfies
\[
-(5c+22)\det M_1
=9w^2(5c+22)-2h^2(32h+2-c).
\]
For \(h\in\mathbb C^\times\), the exact level-one singular-vector
curve is
\begin{equation}\label{eq:w3hol-level-one-null-curve}
9w^2(5c+22)=2h^2(32h+2-c),
\end{equation}
with kernel vector
\[
\left(W_{-1}-\frac{3w}{2h}L_{-1}\right)|h,w\rangle.
\]

\section{The collision-kernel reconstruction}
\label{sec:w3hol-r-matrix-channels}

\begin{theorem}[Four collision channels;
\ClaimStatusConditional]
\label{thm:w3hol-r-channels}
\emph{Type signature:}
\((\mathrm{Open},\mathrm{ordered\ chiral\ bar},2,
H_{W_3}^{\mathrm{bar}}+H_{W_3}^{\mathrm{coll}})\).
With the local coordinate and \(d\log E\) convention fixed by
\(H_{W_3}^{\mathrm{coll}}\), the OPE packet maps to
\begin{align}
K_{TT}^{\mathrm{coll}}(z)
&=\frac{c/2}{z^3}+\frac{2T}{z},
&\operatorname{poles}&=\{3,1\},
\label{eq:w3hol-r-TT}\\
K_{TW}^{\mathrm{coll}}(z)
&=\frac{3W}{z},
&\operatorname{poles}&=\{1\},
\label{eq:w3hol-r-TW}\\
K_{WT}^{\mathrm{coll}}(z)
&=\frac{3W}{z},
&\operatorname{poles}&=\{1\},
\label{eq:w3hol-r-WT}\\
K_{WW}^{\mathrm{coll}}(z)
&=\frac{c/3}{z^5}
+\frac{2T}{z^3}
+\frac{\partial T}{z^2}
+\frac{\frac3{10}\partial^2T
+\frac{32}{5c+22}\Lambda}{z},
&\operatorname{poles}&=\{5,3,2,1\}.
\label{eq:w3hol-r-WW}
\end{align}
The candidate maximal collision order is
\begin{equation}\label{eq:w3hol-kmax}
k_{\max}^{\mathrm{coll}}=5.
\end{equation}
\end{theorem}

\begin{proof}
The exact OPE has pole sets
\(\{4,2,1\}\), \(\{2,1\}\), \(\{2,1\}\), and
\(\{6,4,3,2,1\}\) in the four channels.  The residue convention in
\(H_{W_3}^{\mathrm{coll}}\) pairs the singular product with
\(d\log E\), lowering each retained pole order by one.  Applying this
map term by term gives the displayed conditional kernel.
\end{proof}

\begin{remark}[The even collision pole]
\label{rem:w3hol-even-pole}
The term \(\partial T/(z-w)^3\) in the exact \(WW\) OPE produces the
order-two term \(\partial T/z^2\) in the conditional collision
kernel.  This term records the descendant channel of the weight-three
self-product.
\end{remark}

\section{Modular trace and scalar conductor}
\label{sec:w3hol-kappa}

\begin{theorem}[The genus-one modular coordinate;
\ClaimStatusOpen]
\label{thm:w3hol-kappa-formula}
\emph{Type signature:}
\((\mathrm{Open},\mathrm{modular\ trace},5,
H_{W_3}^{\mathrm{mod}}+H_{W_3}^{\mathrm{res}})\).
The package \(H_{W_3}^{\mathrm{mod}}+H_{W_3}^{\mathrm{res}}\) must
construct a trace for which
\begin{equation}\label{eq:w3hol-kappa-w3}
\operatorname{tr}_1
\operatorname{Obs}^{\mathrm{def}}_{1,1}(\mathcal W_3^k)
=\kappa_{\mathrm{ch}}(\mathcal W_3^k)\lambda_1.
\end{equation}
Computing \(\kappa_{\mathrm{ch}}(\mathcal W_3^k)\) from this trace is
the genus-one \(W_3\) obligation.
\end{theorem}

\begin{remark}[Local norms and modular curvature]
\label{rem:w3hol-channel-decomp}
The exact leading OPE norms are
\begin{equation}\label{eq:w3hol-kappa-channel}
N_T=\frac c2,\qquad N_W=\frac c3.
\end{equation}
Their reciprocal-weight sum is the typed local diagnostic
\(\frac12+\frac13=\frac56\).  The comparison
from \((N_T,N_W)\) to a genus-one curvature belongs to
\(H_{W_3}^{\mathrm{mod}}\).
\end{remark}

\begin{remark}[The anomaly ratio]
\label{rem:w3hol-anomaly-ratio}
On \(c\in\mathbb C^\times\), the modular anomaly ratio is defined by
\begin{equation}\label{eq:w3hol-rho}
\rho_{W_3}^{\mathrm{mod}}(k)
:=\frac{\kappa_{\mathrm{ch}}(\mathcal W_3^k)}{c(k)}.
\end{equation}
Its evaluation is open with the same genus-one trace package.
\end{remark}

\begin{theorem}[The scalar Verdier sum;
\ClaimStatusOpen]
\label{thm:w3hol-kappa-sum}
\emph{Type signature:}
\((\mathrm{Open},\mathrm{bar\text{--}Verdier\ plus\ scalar\ trace},
2\to5,
H_{W_3}^{\mathrm{DS/bar}}+
H_{W_3}^{\mathrm{mod}}+
H_{W_3}^{\mathrm{rescale}})\).
For a constructed Koszul pair, the scalar conductor is
\begin{equation}\label{eq:w3hol-kappa-sum}
K^\kappa(\mathcal W_3^k)
:=
\kappa_{\mathrm{ch}}(\mathcal W_3^k)
+\kappa_{\mathrm{ch}}\bigl((\mathcal W_3^k)^!\bigr).
\end{equation}
Its evaluation requires a family contraction
\((\iota_k,p_k,h_k,\nu_k)\) on one factorization bar--cobar complex,
\[
dh_k+h_kd
=\nu_k(\operatorname{id}-\iota_kp_k),
\]
and a trace map identifying \(\nu_k\) with the displayed scalar.
\end{theorem}

\begin{remark}[Theorem C and the family rescaling problem]
\label{rem:w3hol-theorem-c-scope}
Theorem~C enters after the presentation coalgebra, full bar,
Ran--Verdier algebra, strict Koszul branch, and scalar trace have been
constructed.  For \(W_3\), the proof obligation is the simultaneous
construction of
\[
(\iota_k,p_k,h_k,\nu_k)
\quad\text{and}\quad
\operatorname{Tr}_{W_3}\colon
\mathsf{BarCobar}_{W_3}\longrightarrow
\mathsf{ScalTrace}_{W_3},
\]
with
\(\operatorname{Tr}_{W_3}(\nu_k)=K^\kappa(\mathcal W_3^k)\).
\end{remark}

\section{Stable-graph and scalar-shadow obligations}
\label{sec:w3hol-genus-2}

\begin{theorem}[Genus-two coloured stable-graph sum;
\ClaimStatusOpen]
\label{thm:w3hol-deltaF2}
\emph{Type signature:}
\((\mathrm{Open},\mathrm{modular\ coloured\ graph},5,
H_{W_3}^{\mathrm{bar}}+
H_{W_3}^{\mathrm{proj}}+
H_{W_3}^{\mathrm{sew}})\).
The rank-two genus-two correction is the projected graph sum
\begin{equation}\label{eq:w3hol-deltaF2}
\delta F_{2}^{\mathrm{cross}}(\mathcal W_3^k)
:=
\sum_{\substack{\Gamma\in\mathsf G_{2,0}\\
\sigma:E(\Gamma)\to\{T,W\}\\
\sigma\ \mathrm{uses\ both\ colours}}}
\frac{\operatorname{Cont}_{\Gamma,\sigma}(\Theta_{\mathcal W_3})}
{|\operatorname{Aut}(\Gamma,\sigma)|}.
\end{equation}
The packages in the type signature must construct every vertex,
propagator, projection, and automorphism factor.
\end{theorem}

\begin{equation}\label{eq:w3hol-deltaF2-graph-decomp}
\delta F_{2}^{\mathrm{cross}}
=
\mathsf P_{2}^{W_3}
\left(
\sum_{\Gamma\in\mathsf G_{2,0}}
\operatorname{Cont}_{\Gamma}(\Theta_{\mathcal W_3})
\right).
\end{equation}
This equation fixes the mathematical type of the quantity before its
evaluation.

\begin{proposition}[Genus-three finite-window problem;
\ClaimStatusOpen]
\label{prop:w3hol-deltaF3-finite-window}
\emph{Type signature:}
\((\mathrm{Open},\mathrm{finite\ coloured\ stable\ graphs},5,
H_{W_3}^{\mathrm{bar}}+
H_{W_3}^{\mathrm{proj}}+
H_{W_3}^{\mathrm{sew}})\).
The genus-three cross-channel coordinate is
\begin{equation}\label{eq:w3hol-deltaF3}
\delta F_{3}^{\mathrm{cross}}(\mathcal W_3^k)
:=
\sum_{\substack{\Gamma\in\mathsf G_{3,0}\\
\sigma:E(\Gamma)\to\{T,W\}\\
\sigma\ \mathrm{uses\ both\ colours}}}
\frac{\operatorname{Cont}_{\Gamma,\sigma}(\Theta_{\mathcal W_3})}
{|\operatorname{Aut}(\Gamma,\sigma)|}.
\end{equation}
Under the modular trace package, its diagonal scalar comparison is
\begin{equation}\label{eq:w3hol-genus3-scalar}
F_3^{\mathrm{diag}}(\mathcal W_3^k)
:=
\kappa_{\mathrm{ch}}(\mathcal W_3^k)
\lambda_3^{\mathrm{FP}}.
\end{equation}
\end{proposition}

\begin{corollary}[Criterion for genuinely rank-two genus data;
\ClaimStatusConditional]
\label{cor:w3hol-multi-gen-fails}
\emph{Type signature:}
\((\mathrm{Open},\mathrm{scalar\ plus\ cross\ projection},5,
H_{W_3}^{\mathrm{mod}}+
H_{W_3}^{\mathrm{proj}}+
H_{W_3}^{\mathrm{sew}})\).
If the constructed cross projection represents a nonzero class, then
\begin{equation}\label{eq:w3hol-genus2-rank-two-splitting}
F_2(\mathcal W_3^k)
=F_2^{\mathrm{diag}}(\mathcal W_3^k)
+\delta F_2^{\mathrm{cross}}(\mathcal W_3^k)
\end{equation}
contains information beyond the diagonal scalar coordinate.
\end{corollary}

\begin{remark}[Asymptotic question]
\label{rem:w3hol-large-c}
The large-\(c\) behaviour is defined after
Equation~\eqref{eq:w3hol-deltaF2} has been constructed:
\begin{equation}\label{eq:w3hol-large-c}
\delta F_{2,\infty}^{\mathrm{cross}}
:=
\lim_{c\to\infty}
\delta F_2^{\mathrm{cross}}(\mathcal W_3^k).
\end{equation}
Existence and evaluation of this limit belong to
\(H_{W_3}^{\mathrm{proj}}+H_{W_3}^{\mathrm{sew}}\).
\end{remark}

\begin{remark}[Two weights]
\label{rem:w3hol-uniform-weight-cross-channel}
The weights \(2\) and \(3\) define a two-colour graph problem.
Diagonal and mixed colourings are independent projections of the
same ordered Maurer--Cartan tensor.  This is the first family in the
standard landscape where the mixed projection is a separate
construction.
\end{remark}

\section{Propagator mixing}
\label{sec:w3hol-propagator-variance}

\begin{theorem}[Propagator-mixing coordinate;
\ClaimStatusOpen]
\label{thm:w3hol-propagator-variance}
\emph{Type signature:}
\((\mathrm{Open},\mathrm{two\text{-}channel\ scalar\ projection},5,
H_{W_3}^{\mathrm{res}}+
H_{W_3}^{\mathrm{proj}}+
H_{W_3}^{\mathrm{sew}})\).
Let \(G_{W_3}\) be the projected two-channel propagator metric,
\(f_{W_3}\) the projected quartic gradient, and
\(\mathbf 1=(1,1)\).  The mixing coordinate is
\begin{equation}\label{eq:w3hol-prop-variance}
\delta_{\mathrm{mix}}(\mathcal W_3^k)
:=
\langle f_{W_3},G_{W_3}^{-1}f_{W_3}\rangle
-
\frac{\langle f_{W_3},\mathbf1\rangle^2}
{\langle\mathbf1,G_{W_3}\mathbf1\rangle}.
\end{equation}
The packages in the type signature must construct
\((G_{W_3},f_{W_3})\) from the full Gram blocks and the chain-level
projection.
\end{theorem}

\begin{equation}\label{eq:w3hol-mixing-poly}
P_{\mathrm{mix}}^{W_3}(c)
:=
\det\!\bigl(f_{W_3}(c),G_{W_3}(c)\mathbf1\bigr).
\end{equation}
This determinant records proportionality of the quartic gradient and
the curvature direction once both vectors have been constructed.

\begin{corollary}[Positive-chamber inequality;
\ClaimStatusConditional]
\label{cor:w3hol-prop-variance-positive}
\emph{Type signature:}
\((\mathrm{Open},\mathrm{real\ positive\ propagator\ chamber},5,
H_{W_3}^{\mathrm{proj}}+H_{W_3}^{\mathrm{pos}})\).
If \(G_{W_3}\) is positive definite, then
\[
\delta_{\mathrm{mix}}(\mathcal W_3^k)\geq0.
\]
Equality is equivalent to linear dependence of
\(f_{W_3}\) and \(G_{W_3}\mathbf1\).
\end{corollary}

\section{Scalar-shadow projections}
\label{sec:w3hol-shadow-connection}

\begin{theorem}[The \(T\)-line shadow problem;
\ClaimStatusOpen]
\label{thm:w3hol-Q-T}
\emph{Type signature:}
\((\mathrm{Open},\mathrm{one\text{-}line\ scalar\ projection},5,
H_{W_3}^{\mathrm{res}}+H_{W_3}^{\mathrm{proj}})\).
The projected \(T\)-line metric is
\begin{equation}\label{eq:w3hol-Q-T}
Q_T^{W_3}(t)
:=
\mathsf P_T\!\left(
\operatorname{Sh}(\Theta_{\mathcal W_3});t
\right).
\end{equation}
Its comparison with the Virasoro metric requires an intertwining
theorem for the inclusion generated by \(T\).
\end{theorem}

\begin{theorem}[The \(W\)-line shadow problem;
\ClaimStatusOpen]
\label{thm:w3hol-Q-W}
\emph{Type signature:}
\((\mathrm{Open},\mathrm{one\text{-}line\ scalar\ projection},5,
H_{W_3}^{\mathrm{res}}+H_{W_3}^{\mathrm{proj}})\).
The projected \(W\)-line metric is
\begin{equation}\label{eq:w3hol-Q-W}
Q_W^{W_3}(t)
:=
\mathsf P_W\!\left(
\operatorname{Sh}(\Theta_{\mathcal W_3});t
\right).
\end{equation}
The involution \(W\mapsto-W\) supplies its parity grading.
\end{theorem}

\begin{theorem}[Projected discriminants;
\ClaimStatusOpen]
\label{thm:w3hol-discriminants}
\emph{Type signature:}
\((\mathrm{Open},\mathrm{one\text{-}line\ scalar\ projection},5,
H_{W_3}^{\mathrm{res}}+H_{W_3}^{\mathrm{proj}})\).
After the metrics in
Theorems~\ref{thm:w3hol-Q-T} and~\ref{thm:w3hol-Q-W} have been
constructed, define
\begin{equation}\label{eq:w3hol-discriminants}
\Delta_T^{W_3}:=\operatorname{Disc}_t Q_T^{W_3}(t),
\qquad
\Delta_W^{W_3}:=\operatorname{Disc}_t Q_W^{W_3}(t).
\end{equation}
Their evaluation and shadow-depth interpretation are finite
projection problems.
\end{theorem}

\begin{remark}[Exact exchange data and projected contacts]
\label{rem:w3hol-S4}
The exact rank-one \(\Lambda\)-exchange ratios are
\[
\frac{(32/D(c))^2}{N_\Lambda}
=\frac{10240}{cD(c)^3},
\qquad
\frac{(16/D(c))^2}{N_\Lambda}
=\frac{2560}{cD(c)^3}.
\]
They use OPE and mode normalizations, respectively.  The projected
quartic contacts are defined by
\begin{equation}\label{eq:w3hol-S4}
S_{4,T}^{W_3}:=\mathsf P_T^{(4)}(\Theta_{\mathcal W_3}),
\qquad
S_{4,W}^{W_3}:=\mathsf P_W^{(4)}(\Theta_{\mathcal W_3}),
\end{equation}
under \(H_{W_3}^{\mathrm{res}}+H_{W_3}^{\mathrm{proj}}\).
\end{remark}

\section{Hamiltonians and the holographic lift}
\label{sec:w3hol-commuting-hamiltonians}

\begin{theorem}[Commuting Hamiltonians from the represented
connection; \ClaimStatusConditional]
\label{thm:w3hol-commuting-differentials}
\emph{Type signature:}
\((\mathrm{Open}\times\mathrm{CY},
\mathrm{represented\ ordered\ bar},4,
H_{W_3}^{\mathrm{coll}}+
H_{W_3}^{\mathrm{flat}}+
H_{W_3}^{\mathrm{OC}})\).
Let
\[
H_i^{W_3}
:=
\rho_{\mathrm{conf}}
\bigl(
\operatorname{Sh}_{0,n}(\Theta_{\mathcal W_3})
\bigr)_i.
\]
If \(H_{W_3}^{\mathrm{flat}}\) identifies the Maurer--Cartan curvature
with the curvature of this representation, then
\begin{equation}\label{eq:w3hol-commuting}
[H_i^{W_3},H_j^{W_3}]=0
\qquad(i\ne j).
\end{equation}
Under the collision-order comparison, the candidate differential
order is bounded by \(k_{\max}^{\mathrm{coll}}-1=4\).
\end{theorem}

\begin{remark}[Sphere reconstruction]
\label{rem:w3hol-gz26}
A comparison with a Toda or Gaudin differential system consists of a
map between conformal-block bundles, a coordinate identification, and
a normalization equality for the binary kernels.  These data form a
specialization of
\(H_{W_3}^{\mathrm{coll}}+H_{W_3}^{\mathrm{flat}}\).
\end{remark}

The seven entries assemble under \(H_{W_3}^{\mathrm{hol}}\):
\[
\mathcal H(\mathcal W_3^k)
=
\left(
\mathcal W_3^k,\,
(\mathcal W_3^k)^{\mathrm i},\,
(\mathcal W_3^k)^!,\,
Z_{\mathrm{ch}}^{\mathrm{der}}(\mathcal W_3^k),\,
K_{\mathcal W_3}^{\mathrm{coll}},\,
\Theta_{\mathcal W_3},\,
\nabla^{\mathrm{hol}}
\right).
\]
The exact local algebra supplies the first entry and the finite
constraints on the remaining six.  The named packages supply the
state transitions which turn those constraints into objects and
comparison maps.

\section{Computational surface}
\label{sec:w3hol-verification}

\begin{center}
\renewcommand{\arraystretch}{1.25}
\begin{tabular}{lll}
\toprule
\textbf{Datum} & \textbf{Status} & \textbf{Anchor}\\
\midrule
\(c(k)\) & exact & Fateev--Lukyanov DS formula\\
\(c(k)+c(-k-6)=100\) & exact & rational-function arithmetic\\
\(WW\to\Lambda,\partial\Lambda\) & exact & \(32/D(c),16/D(c)\)\\
\(N_\Lambda\) & exact & \(cD(c)/10\)\\
\(\Lambda_0|h,w\rangle\) & exact & \((h^2+h/5)|h,w\rangle\)\\
\(\det M_1\) & exact & Equation~\eqref{eq:w3hol-level-one-null-curve}\\
\(\mathcal A^{\mathrm i}\to\bar B^{\mathrm{ch}}(\mathcal A)\)
& open & \(H_{W_3}^{\mathrm{pres}}\)\\
\(\mathcal A^{\mathord{!}}\) & conditional
& \(H_{W_3}^{\mathrm{DS/bar}}\)\\
\(Z_{\mathrm{ch}}^{\mathrm{der}}(\mathcal A)\) comparison
& open & \(H_{W_3}^{\mathrm{cen}}+H_{W_3}^{\mathrm{OC}}\)\\
\(K_{\mathcal A}^{\mathrm{coll}}\) & conditional
& \(H_{W_3}^{\mathrm{bar}}+H_{W_3}^{\mathrm{coll}}\)\\
\(\kappa_{\mathrm{ch}},\rho_{W_3}^{\mathrm{mod}},K^\kappa\)
& open & \(H_{W_3}^{\mathrm{mod}}+H_{W_3}^{\mathrm{rescale}}\)\\
stable-graph sums & open
& \(H_{W_3}^{\mathrm{proj}}+H_{W_3}^{\mathrm{sew}}\)\\
scalar shadows & open
& \(H_{W_3}^{\mathrm{res}}+H_{W_3}^{\mathrm{proj}}\)\\
Hamiltonians & conditional
& \(H_{W_3}^{\mathrm{flat}}+H_{W_3}^{\mathrm{OC}}\)\\
\bottomrule
\end{tabular}
\end{center}

The companion engine
\texttt{compute/lib/theorem\_w3\_holographic\_datum\_engine.py}
checks the exact rows symbolically and represents every reconstruction
row by a typed open or conditional packet.  Its tests verify the
central reflection, the \(32/16\) OPE normalization, \(N_\Lambda\),
the \(\Lambda_0\) action, the level-one determinant, and every
hypothesis boundary.
